\section{Qualitative Predictions}


The framework of this paper is not a theory that measures consciousness as a single quantity.

The predictions presented here are therefore not numerical predictions of a "consciousness quantity" but qualitative predictions about how the conscious state deforms when speed difference, coupling viscosity, internal delay, meaning gradient, and effective torque-like quantity change.

The central relation is the effective quantity introduced in preceding sections:

\begin{equation}
\mathcal{T}_{\Phi} \sim \kappa \frac{\Delta v}{\mu} \left|\nabla \Phi\right|
\end{equation}

$\mathcal{T}_{\Phi}$ is the effective torque-like quantity coupling time-scale difference and meaning gradient. It does not generate consciousness; it is a structural index showing how strongly and in which direction log referencing tends to open.

\subsection{Synchronization Collapse}


This prediction concerns the case where speed difference is too small:

\begin{equation}
\Delta v \rightarrow 0
\end{equation}

The boundary between the fast and slow layers collapses into synchrony. Processing stabilizes, but the temporal thickness required for referencing is lost. Conscious access weakens; action tends toward habituation, automation, and reflexive response.

In this state, processing occurs, but the sense of "currently referencing" diminishes.

Excessive synchrony therefore does not strengthen consciousness; it works toward making conscious referencing unnecessary.

\subsection{Dissociation by Excessive Speed Differential}


This prediction concerns the case where speed difference is too large:

\begin{equation}
|\Delta v| \gg 0
\end{equation}

The correspondence between fast and slow layers becomes difficult to maintain. Predictions and associations run in the fast layer, but they do not stably sediment as logs in the slow layer. Conscious access fragments; structures resembling dreaming, daydreaming, dissociation, and hallucinatory states become more likely to appear.

In this state, even though processing intensity is high, the coherence of log referencing decreases. Referencing does not strengthen — it scatters.

\subsection{Coupling Viscosity and State Transition}


This prediction concerns changes in coupling viscosity $\mu$.

When $\mu$ is too low, speed difference does not slip sufficiently and the layers tend to adhere:

\begin{equation}
\mu \rightarrow 0
\end{equation}

Processing becomes rigid and tends to fix to specific logs or meaning wells.

When $\mu$ is too high, the layers slip excessively:

\begin{equation}
\mu \rightarrow \infty
\end{equation}

Fast-layer changes do not transfer sufficiently to the slow layer; log referencing becomes unstable.

The conscious state is therefore most stable in the intermediate range of $\mu$:

\begin{equation}
\mu_{\min} < \mu < \mu_{\max}
\end{equation}

This prediction suggests the possibility of treating conscious states as a phase diagram: consciousness is not binary ON/OFF but a regime that transitions continuously with coupling viscosity and speed difference.

\subsection{Delay Window}


This prediction concerns internal delay $\tau$:

\begin{equation}
\tau \propto \frac{\Delta v}{\mu}
\end{equation}

When $\tau$ is too small, processing approaches immediate reaction and disappears before sedimenting into the slow layer ($\tau < \tau_{\min}$). Conscious referencing is weak and action approaches automatic processing.

When $\tau$ is too large, the correspondence between current processing and log sedimentation breaks down ($\tau > \tau_{\max}$). Consciousness lags, fragments, and instabilities in self-sense and time sense become more likely.

A stable delay window is therefore required:

\begin{equation}
\tau_{\min} < \tau < \tau_{\max}
\end{equation}

\subsection{Meaning Gradient Bias}


This prediction concerns meaning gradient $\nabla\Phi$.

Log referencing is not opened equally for all logs. In regions where the gradient of the information potential $\Phi$ is large, processing is easily drawn in that direction:

\begin{equation}
|\nabla \Phi| \uparrow \quad \Rightarrow \quad \text{referencing bias} \uparrow
\end{equation}

Memories, emotions, concepts, and narratives with strong meaning gradients more easily enter conscious referencing.

But when the meaning gradient is too strong, referencing becomes fixed to specific wells. This can appear as repetitive thought, compulsive attention, fixation, and traumatic re-referencing.

The meaning gradient directs consciousness, but an excessive gradient takes away the degrees of freedom of referencing.

\subsection{Qualia and Measurement}


This prediction concerns qualia and measurement.

If qualia are boundary phenomena between fast and slow layers, measurement changes that boundary.

Measurement converts live boundary states into slow, reportable logs:

\begin{equation}
\text{measurement} \quad \Rightarrow \quad \text{boundary closure} \quad \Rightarrow \quad \text{sedimented residue}
\end{equation}

This prediction is not the simple claim that qualia cannot be measured.

It is the claim that what can be measured is not qualia themselves but the log structure remaining after qualia have passed through.

\subsection{Measurement-Induced Residue Shift}


The stronger the measurement demand, the more live qualia sediment into the slow layer and stabilize as reportable residue.

Increased measurement precision therefore appears not as capture of qualia themselves but as clarification of sedimented residue.

This prediction requires treating measurement not as mere observation but as a regime-changing operation on the conscious state.

Even for the same experience, if the strength of reporting demand, options, verbalization, or numerical evaluation changes, the remaining log structure changes.

Qualia reports obtained in experiments must therefore be interpreted not as live boundary states themselves but as residue shaped by measurement conditions.

\subsection{Split-Brain Asymmetry}


This prediction concerns split-brain experiments.

Consciousness does not double in split-brain patients.

What is predicted is the asymmetrization of log-referencing pathways.

Even if processing occurs in one pathway, if it cannot reach a slow reportable log, it does not enter explicit consciousness.

The rationalization observed in split-brain patients is therefore not another self speaking but the slow layer supplying coherence for inaccessible log gaps.

\subsection{Non-Local Regime Prediction}


If conscious states are temporal regimes rather than single sites, then even the same local activity will produce different conscious reports when speed difference, delay, meaning gradient, and reporting conditions change.

Neural correlates cannot alone determine conscious states and must be interpreted together with temporal coupling conditions.

This prediction does not deny local correspondence of neural activity.

Rather, it claims that the same site activity carries different conscious meaning depending on which temporal coupling it participates in within the consciousness window.

Therefore, even with identical site activity, if slow-log sedimentation, connection to reporting pathways, and meaning gradient strength differ, subjective report and behavioral response should differ.

\subsection{Consciousness Window as Phase Diagram}


The predictions above show that conscious states exist not as a point but as a window.

\begin{equation}
0 < |\Delta v| < \Delta v_{\max}, \quad \mu_{\min} < \mu < \mu_{\max}, \quad \tau_{\min} < \tau < \tau_{\max}, \quad \mathcal{T}_{\min} < \mathcal{T}_{\Phi} < \mathcal{T}_{\max}
\end{equation}

Within this consciousness window, processing does not fully synchronize or fully separate, and log referencing opens stably.

This window can be read as a phase diagram. Placing speed difference $|\Delta v|$ on the horizontal axis and coupling viscosity $\mu$ on the vertical axis, synchronization collapse, dissociation, fixation, fragmentation, and stable log referencing are positioned as distinct regimes.

Internal delay $\tau$ and meaning gradient $|\nabla\Phi|$ function as control variables that extend the state on this phase diagram in a depth direction.

The consciousness window is therefore not an ON/OFF boundary but a finite stable region formed by multiple control variables.

Outside the window, conscious states collapse in different directions:

\begin{itemize}
\item Collapse toward synchrony: automation, habituation, and reflexivity strengthen.
\item Collapse toward separation: fragmentation, dream-like states, dissociation, and hallucinatory states strengthen.
\item Too-low viscosity: fixation to meaning wells strengthens.
\item Too-high viscosity: continuity of referencing is lost.
\item Too-weak meaning gradient: referencing loses direction.
\item Too-strong meaning gradient: referencing becomes fixed.
\end{itemize}

\subsection{Testable Directions}


This paper does not treat consciousness as a single directly measurable quantity.

However, directions for verifying conscious state transitions from changes in peripheral quantities are open.

Among verifiable directions, at minimum:

\begin{enumerate}
\item Correspondence between the time-scale difference of neural activity and the stability of subjective report.
\item Changes in $\tau$ through attentional manipulation.
\item Increase in referencing bias for semantically loaded stimuli.
\item Transformation of qualia-like reports by strengthened measurement and report demands.
\item Changes in speed difference, viscosity, and delay window in sleep, dreaming, anesthesia, and dissociative states.
\item Asymmetry of log-referencing pathways in linguistically non-reportable processing.
\item Clarification of residue and transformation of live qualia reports through measurement demand intensity.
\item Changes in conscious report due to differences in temporal coupling conditions for identical local activity.
\item Phase-diagram classification of conscious states on axes of speed difference and coupling viscosity.
\end{enumerate}

None of these complete the framework of this paper.

Rather, they are entry points for investigating consciousness not as a product to be found but as a non-closure state across time scales.

\subsection{Summary}


The qualitative predictions of this paper reduce to one point:

Conscious states stabilize when speed difference, coupling viscosity, internal delay, and meaning gradient enter an intermediate effective region:

\begin{equation}
\text{stable consciousness} \quad \Leftrightarrow \quad (\Delta v, \mu, \tau, \mathcal{T}_{\Phi}, \nabla\Phi) \in \text{consciousness window}
\end{equation}

This window is not a fixed location.

It is the dynamic region in which temporal non-closure is maintained without rupture.

Consciousness is therefore not modeled here as a generated object.

Consciousness opens when conditions are met, and closes when conditions break.
